% Einheit 3 von 3 – Waagerechte Asymptoten (bank/grenzwerte-und-verhalten-im-unendlichen/e3.jsonl, zusammenbau v0.1)
%% TODO Zeitmarke Einheit 3: Marken-Zeile nicht lesbar
%% TODO Zweigzeile Teil 1: was hier gelernt wird (Satz in Schülersprache aus dem Einheitstitel) – Einheit 3
\einheitenkopf[e3]{Einheit 3 von 3 · Waagerechte Asymptoten}
\zweigzeile{keine Prüfungsaufgabe}
\setcounter{aufgabe}{14}

%% TODO Titel: Ich-kann-Satz fehlt in der Bank (Platzhalter: Kettenname) – Nr. 15
%% TODO Anweisung: ein Satz über der ersten Teilaufgabe fehlt in der Bank – Nr. 15
\begin{aufgabe}{Waagerechte Asymptoten}
%% TODO \rechenplatz{n}: Schrittzahl des Grundfalls fehlt in der Bank (nur bei fester Schreibform, 2.3 b)
\begin{teile}
\teil Wer gewinnt für $x \to +\infty$? $f(x) = 2e^{-x} - 7$ \kreuz{der Summand $2e^{-x}$} \\ \kreuz{der Summand $-7$}
\teil $f(x) = 2e^{-x} + 4$. Grenzwert für $x \to +\infty$ und waagerechte Asymptote? f(x) → \leerfeld ; Asymptote y = \leerfeld
\teil $f(x) = 4e^{-0{,}5x} - 4$. Begründe, dass $f$ streng monoton fällt und der Graph durch den Ursprung verläuft. Grenzwert für $x \to +\infty$? \hfill (Abitur 2026 LK)
\teil $f(x) = \frac{8}{1 + e^{x}}$. Begründe, dass $f$ keine Nullstelle hat, und gib die Grenzwerte für $x \to \pm\infty$ an. \hfill (Abitur 2024 LK)
\teil Ein Getränk kühlt im Kühlschrank ab: $T(t) = 18e^{-0{,}1t} + 4$ ($t \ge 0$ in min, $T$ in °C). Gib den Wertebereich an und deute den Grenzwert im Sachzusammenhang.
\end{teile}
\end{aufgabe}

%% TODO Titel: Ich-kann-Satz fehlt in der Bank (Platzhalter: Kettenname) – Nr. 16
%% TODO Anweisung: ein Satz über der ersten Teilaufgabe fehlt in der Bank – Nr. 16
\begin{aufgabe}{Waagerechte Asymptoten – Fehler finden, Begründen, Anwendung}
\begin{teile}
\teil $f(x) = 2e^{-x} - 9$. Ein Schüler schreibt: „$f(x) \to 0$ für $x \to +\infty$, weil $e^{-x}$ gegen null geht.“ Finde den Fehler.
\teil Begründe, warum der Graph von $f(x) = 5e^{-x} + 3$ für große $x$ der Geraden $y = 3$ beliebig nahe kommt.
\teil Ein Kuchen kühlt nach $T(t) = 20 + 160e^{-0{,}05t}$ ab ($t$ in min, $T$ in °C). Welche Temperatur erreicht er langfristig, und warum ist das sinnvoll?
\end{teile}
\end{aufgabe}
